\section{Prediction Workflow}

\subsection{Random Forest Pipeline}
The pipeline uses the cleaned churn dataset,
creates grouped features, encodes categorical variables, trains a Random Forest
classifier, and exports model outputs for later analysis.

\subsection{Model Evaluation Summary}
Table \ref{tab:test_metrics} reports the Random Forest test metrics:
\begin{table}[h]
\centering
\caption{Random Forest test metrics from the group repository}
\label{tab:test_metrics}
\begin{tabular}{lr}
\toprule
\textbf{Metric} & \textbf{Value} \\
\midrule
ROC AUC & 0.676510 \\
Accuracy & 0.761694 \\
Precision & 0.187745 \\
Recall & 0.421945 \\
F1-score & 0.259863 \\
\bottomrule
\end{tabular}
\end{table}

\subsection{Model Performance Interpretation}
The model performance should be interpreted in the context of churn prediction,
rather than as a generic classification score. In the working dataset, the
number of churned customers is 83,362 out of 840,765 customers, which gives a
baseline churn rate of approximately 9.92\%. This means that the target class is
strongly imbalanced: only about 10 out of every 100 customers are churned
customers.

For this reason, accuracy is not the primary evaluation metric. A naive model
that predicts every customer as non-churn would already obtain about 90.08\%
accuracy, but it would detect no churned customers. The Random Forest accuracy
of 76.17\% should therefore not be interpreted as the main indicator of model
quality. More informative metrics are precision, recall, F1-score, ROC AUC, and
precision-recall based metrics. In this run, the recall of 42.19\% means that
the model identifies about 42\% of the actual churned customers, while the low
precision reflects the difficulty of separating a relatively small churn class
from a much larger non-churn class.

The model is also constrained by the Phase 2 feature design. The initial model
uses interpretable, dashboard-friendly grouped features such as contract type,
customer satisfaction group, complaint group, service-call group, late-payment
group, monthly-charge group, tenure group, service-count group, and technical
support. These variables contain useful directional signals, but they do not
capture all possible churn drivers. For example, the dataset does not include
detailed behavioral histories such as plan changes, discount history, competitor
offers, support ticket text, customer sentiment, or detailed payment timelines.

Another limitation is that several numeric variables are grouped into bins before
modeling. This helps the dashboard remain readable and supports clear
segment-level interpretation, but it can remove continuous detail that a
predictive model might otherwise use. Therefore, the Random Forest model should
be treated as a decision-support component for prioritizing retention actions,
not as a perfect churn prediction system.

The selected threshold is chosen on the validation set to improve the balance
between precision and recall. Lowering the threshold would identify more churned
customers but would also increase false positives. Raising the threshold would
improve precision but miss more churned customers. This trade-off is acceptable
for a retention dashboard, where the objective is to highlight customers or
segments that deserve earlier attention from customer support or retention
teams.

Future model improvement should focus on using raw numeric features together
with grouped features, testing gradient boosting models such as XGBoost,
LightGBM, or CatBoost, tuning hyperparameters against F1 or PR-AUC rather than
accuracy, handling class imbalance with class weights or resampling inside the
training pipeline, and adding richer behavioral features if they become
available.
